\documentclass[tikz,border=6pt]{standalone}
\usepackage[no-math]{fontspec}
\usepackage[brazil]{babel}
\usepackage{amsmath,amssymb}
\IfFileExists{C:/Windows/Fonts/Lato_Regular.ttf}{%
  \setmainfont{Lato}[
    Path = C:/Windows/Fonts/,
    UprightFont = Lato_Regular.ttf,
    BoldFont = Lato_Bold.ttf,
    ItalicFont = Lato_Italic.ttf,
    BoldItalicFont = Lato_BoldItalic.ttf
  ]
  \setsansfont{Lato}[
    Path = C:/Windows/Fonts/,
    UprightFont = Lato_Regular.ttf,
    BoldFont = Lato_Bold.ttf,
    ItalicFont = Lato_Italic.ttf,
    BoldItalicFont = Lato_BoldItalic.ttf
  ]
}{%
  \RequirePackage[default]{lato}
}
\renewcommand{\familydefault}{\sfdefault}

\usepackage{xcolor}
\usepackage{tikz}
\usepackage{pgfplots}
\pgfplotsset{compat=1.18}
\usetikzlibrary{arrows.meta,positioning,shapes.geometric,fit,backgrounds,calc,decorations.pathreplacing,patterns}

% Paleta institucional UEFS / CHF662
\definecolor{uefsNavy}{HTML}{1B2A4A}
\definecolor{uefsTeal}{HTML}{0E7C7B}
\definecolor{uefsForest}{HTML}{2D6A4F}
\definecolor{uefsGold}{HTML}{C4983B}
\definecolor{uefsGoldDark}{HTML}{8F6B1E}
\definecolor{uefsAlert}{HTML}{B02A37}
\definecolor{uefsOrange}{HTML}{D95F02}
\definecolor{uefsInk}{HTML}{22252A}
\definecolor{uefsMuted}{HTML}{5C6470}
\definecolor{uefsWater}{HTML}{2E86AB}
\definecolor{uefsSand}{HTML}{D9B382}
\definecolor{bgNavy}{HTML}{F2F5F9}
\definecolor{bgTeal}{HTML}{E9F5F4}
\definecolor{bgForest}{HTML}{EEF5F0}
\definecolor{bgGold}{HTML}{FCF7EB}
\definecolor{bgAlert}{HTML}{FDF2F2}
\definecolor{bgGray}{HTML}{F8FAFC}
\definecolor{lineGray}{HTML}{CFD8DC}

\tikzset{
  >=Stealth,
  aaCard/.style={draw=#1, fill=white, rounded corners=4pt, line width=1pt, inner sep=6pt, align=center, font=\small},
  aaBox/.style={draw=#1, rounded corners=4pt, line width=1.1pt, inner sep=6pt, align=center, font=\footnotesize\sffamily, text=uefsInk},
  aaTitle/.style={font=\normalsize\bfseries\sffamily, text=uefsNavy, align=center},
  aaSub/.style={font=\footnotesize\sffamily, text=uefsMuted, align=center},
  aaArrow/.style={->, >=Stealth, line width=1.4pt},
  aaDArrow/.style={<->, >=Stealth, line width=1.4pt}
}


\begin{document}
\begin{tikzpicture}[font=\sffamily]

% =========================================================
% PAINEL (a) Fluxo Metodológico da AIA (6 Etapas)
% =========================================================
\node[aaTitle] at (-4.8, 5.25) {(a) Fluxo Metodológico da Avaliação de Impacto Ambiental};

\begin{scope}[shift={(-4.8, 0)}]
  \node[aaBox=uefsNavy, fill=bgNavy, text width=8.6cm, minimum height=0.95cm, inner sep=6pt] (s1) at (0, 4.1)
    {\textbf{1. Triagem e Definição de Escopo (Termo de Referência)}};

  \node[aaBox=uefsTeal, fill=bgTeal, text width=8.6cm, minimum height=0.95cm, inner sep=6pt] (s2) at (0, 2.55)
    {\textbf{2. Diagnóstico Ambiental (Linha de Base: Físico, Biótico, Socioeconômico)}};

  \node[aaBox=uefsGold, fill=bgGold, text width=8.6cm, minimum height=0.95cm, inner sep=6pt] (s3) at (0, 1.0)
    {\textbf{3. Identificação, Previsão e Valoração dos Impactos (Matrizes / Redes)}};

  \node[aaBox=uefsNavy, fill=bgNavy, text width=8.6cm, minimum height=0.95cm, inner sep=6pt] (s4) at (0, -0.55)
    {\textbf{4. Prognóstico (Cenário com Projeto vs. Sem Projeto)}};

  \node[aaBox=uefsForest, fill=bgForest, text width=8.6cm, minimum height=0.95cm, inner sep=6pt] (s5) at (0, -2.1)
    {\textbf{5. Medidas Mitigadoras, Compensatórias e Planos de Gestão (PBA)}};

  \node[aaBox=uefsAlert, fill=bgAlert, text width=8.6cm, minimum height=0.95cm, inner sep=6pt] (s6) at (0, -3.65)
    {\textbf{6. Decisão Pública, Monitoramento e Auditoria Pós-Implantação}};

  \draw[aaArrow, uefsMuted, shorten >=2pt, shorten <=2pt] (s1.south) -- (s2.north);
  \draw[aaArrow, uefsMuted, shorten >=2pt, shorten <=2pt] (s2.south) -- (s3.north);
  \draw[aaArrow, uefsMuted, shorten >=2pt, shorten <=2pt] (s3.south) -- (s4.north);
  \draw[aaArrow, uefsMuted, shorten >=2pt, shorten <=2pt] (s4.south) -- (s5.north);
  \draw[aaArrow, uefsMuted, shorten >=2pt, shorten <=2pt] (s5.south) -- (s6.north);
\end{scope}

% =========================================================
% PAINEL (b) Hierarquia de Mitigação (Ordem de Prioridade)
% =========================================================
\node[aaTitle] at (5.3, 5.25) {(b) Hierarquia de Mitigação (Ordem de Prioridade)};

\begin{scope}[shift={(5.3, 0)}]
  % Degraus em pirâmide invertida, com largura suficiente para o texto
  \node[draw=uefsForest, fill=uefsForest, rounded corners=4pt, text=white,
        align=center, font=\footnotesize\sffamily, text width=7.2cm, minimum height=1.55cm, inner sep=6pt]
    (h1) at (0.65, 3.6)
    {\textbf{1º EVITAR (Prevenção na origem)}\\[2pt]Alternativa locacional ou tecnológica};

  \node[draw=uefsTeal, fill=uefsTeal, rounded corners=4pt, text=white,
        align=center, font=\footnotesize\sffamily, text width=6.4cm, minimum height=1.55cm, inner sep=6pt]
    (h2) at (0.65, 1.55)
    {\textbf{2º MINIMIZAR (Redução de magnitude)}\\[2pt]Controle de erosão, filtros, cronograma};

  \node[draw=uefsGoldDark, fill=uefsGold!90!uefsGoldDark, rounded corners=4pt, text=white,
        align=center, font=\footnotesize\sffamily, text width=5.6cm, minimum height=1.55cm, inner sep=6pt]
    (h3) at (0.65, -0.5)
    {\textbf{3º RECUPERAR / REABILITAR}\\[2pt]PRAD, revegetação, bioengenharia};

  \node[draw=uefsAlert, fill=uefsAlert, rounded corners=4pt, text=white,
        align=center, font=\footnotesize\sffamily, text width=4.8cm, minimum height=1.55cm, inner sep=6pt]
    (h4) at (0.65, -2.55)
    {\textbf{4º COMPENSAR (Impacto residual)}\\[2pt]Unidades de conservação, reposição};

  % Seta lateral limpa de preferência ecológica (totalmente separada dos blocos)
  \draw[->, >=Stealth, line width=2pt, uefsForest] (-3.75, -2.4) -- (-3.75, 4.1);
  \node[align=center, font=\scriptsize\bfseries, text=uefsForest, text width=2.2cm]
    at (-3.75, -3.25) {Preferível\\(menor custo\\ecológico)};
\end{scope}

\end{tikzpicture}
\end{document}
